\songtitle{No te pido una respuesta}

\tag*{2019}\tag{hip-hop}\tag{spanish-lyrics}\tag{original-source}
% -> internal.tex (Essays?)
\begin{notes}
Lyrics by Suno based on part of Carlos Thompson's blog entry \emph{Inclusivismos y fundamentalismos} (November 2019): the unknown neighbor next door and what makes him uncomfortable before judging. Retrieved from Wyomee Dank's Suno page (V6-MINI, 15 November 2026).
\end{notes}
\begin{style}
Colombian hip hop musical theater, Spanish-language theatrical lead vocal with clear central Colombian diction, rhythmically nimble spoken-sung verses and a memorable melodic ensemble refrain. Start with sparse piano and close voice, build into syncopated boom-bap drums, warm bass, hand percussion, and small brass accents; shift between intimate questioning and layered call-and-response chorus, then strip back for a reflective ending. Clear, organic stage-like production with punchy drums and intelligible vocals.
\end{style}
\begin{lyrics}
\songpart{Intro}
En la puerta de al lado\\
vive alguien que no conozco.\\
Antes de juzgarlo,\\
me pregunto qué me incomoda.

\songpart{Verse 1}
¿Qué tan distinto puede ser mi vecino,\\
el del otro lado, el del mismo corredor?\\
¿Puedo aceptar que tenga otra piel,\\
que venga de Venezuela o de Estados Unidos?\\
¿Y si ama a un hombre, o si es una pareja?\\
¿Y si su nombre y su cuerpo no encajan en mi idea?\\
¿Si viene de otro estrato, habla distinto al mío,\\
y usa palabras que yo nunca aprendí de niño?

\annotation{Refrain}\\
No tienes que responderme,\\
no te estoy pasando lista.\\
Yo también tengo mis límites,\\
mis costumbres, mis heridas.\\
Pero mira dónde los pones,\\
qué temor los dibujó:\\
¿es la vida de esa persona,\\
o el espejo que me mostró?

\songpart{Verse 2}
¿Puedo vivir junto a un cristiano fundamentalista,\\
un musulmán devoto, un ateo convencido?\\
¿Con el que vota por Uribe, con quien vota por Petro,\\
con marxistas, liberales, con el centro en movimiento?\\
¿Con el costeño que conversa con todo el edificio,\\
con el político corrupto que se esconde en el saludo?\\
¿Con un líder social, con quien fuma marihuana?\\
Y hay preguntas que no son lo mismo, que exigen otra balanza:\\
¿cómo pongo límites al daño y protejo a quien está en riesgo,\\
sin usar una diferencia como excusa para el desprecio?

\annotation{Refrain}\\
No tienes que responderme,\\
no te estoy pasando lista.\\
Yo también tengo mis límites,\\
mis costumbres, mis heridas.\\
Pero mira dónde los pones,\\
qué temor los dibujó:\\
¿es la vida de esa persona,\\
o el espejo que me mostró?

\songpart{Bridge}
No es lo mismo una opinión\\
que una amenaza o una agresión.\\
No es lo mismo poner cuidado\\
que cerrar el mundo con candado.\\
En mi barrio, en mi pueblo,\\
en la calle y la ciudad,\\
¿a quién no quisiera encontrar?\\
¿Y por qué? ¿Y por qué será?

\songpart{Verse 3}
Si mis hijos, cuando crezcan, eligen otro camino,\\
¿qué decisión me haría sentir que fracasé como destino?\\
¿Que sean ateos, comunistas, que amen a quien quieran,\\
o que vivan sin pedirme permiso para ser quienes eran?\\
¿Me dolería su elección o que no se parezcan a mí?\\
¿Confundiría mi miedo con quererlos proteger?\\
La pregunta no se acaba cuando cierro la ventana:\\
¿quién cabe en mi ciudad, quién merece estar mañana?

\songpart{Final Chorus}
No tienes que responderme,\\
solo quédate un momento.\\
Yo también tengo mis límites,\\
pero puedo examinarlos por dentro.\\
En la casa de al lado,\\
en la cuadra, en la ciudad,\\
¿qué diferencia me asusta
 y qué daño debo parar?

\songpart{Outro}
No te pido una respuesta.\\
Empieza por preguntar:\\
¿qué límites son cuidado,\\
y cuáles no dejan pasar?
\end{lyrics}

